\section{Background}\label{sec:background}

This section grounds cA2A in the protocols and mechanisms it composes. We describe A2A
and its Signed Agent Card, the tool-side analog in MCP and cMCP, capability attenuation
and how agent-manifest realizes a signed delegation chain, trusted execution environments
and remote attestation, the Cedar policy language, and the provenance and transparency
work that the TRACE record format draws on. None of this section is a novel claim; it
fixes terminology and states what each cited system provides.

\subsection{A2A and the Signed Agent Card}

The Agent2Agent (A2A) protocol~\citep{a2a2026} is a protocol for agent-to-agent
communication. An agent publishes an \emph{Agent Card}, a JSON document
that advertises the agent's identity, endpoint, and the skills it offers. A client agent
reads a peer's card, then sends it a \emph{task}, a unit of work the peer executes and
reports back on. A2A defines how tasks are moved and how their state transitions; it does
not define an authorization model for what a peer is permitted to do with a task once it
holds it.

The Signed Agent Card is A2A's trust anchor. The card carries a signature that lets a
client verify that the claimed domain owner issued it. This answers one question: did this
domain publish this card. It does not establish that the peer is running attested code,
that a delegating agent holds the authority it forwards, that the task payload is
confidential to the peer, or that there is an unbroken record of who delegated what to
whom. cA2A adds these guarantees as a profile on A2A rather than as a competing transport.

\subsection{MCP and cMCP as the Tool-Side Analog}

The Model Context Protocol (MCP)~\citep{mcp2024} standardizes how an agent discovers and
invokes tools hosted on external servers. Where A2A governs agent-to-agent calls, MCP
governs the agent-to-tool call. MCP, like A2A, defines discovery and invocation but not
authorization: it specifies no policy model, requires no audit record, and gives a client
no way to confirm that a cached tool definition is still the definition the server will
execute.

cMCP~\citep{cmcp2026} closes those gaps for the tool call. It interposes a gateway that
evaluates a Cedar policy per call, maintains a tamper-evident audit chain, and emits a
signed TRACE Claim that any verifier can check offline. cA2A is the agent-to-agent
counterpart of the same program: it takes the tool-side pattern of attested, policy-bound,
provable calls and applies it to the delegation edge between agents. cA2A reuses cMCP's
TEE provider pattern and runs the same Cedar engine (\texttt{cedarpy}) for local policy.

\subsection{Capability Attenuation and Macaroons}

The confused deputy problem~\citep{confuseddeputy1988} arises when a component with
authority is induced to exercise it on behalf of a caller who should not have that
authority. In a multi-agent workflow the deputy is any agent that accepts a narrow task
and then acts with broader authority than the task granted. The capability discipline
answers this by making authority a value that is passed explicitly and can only be
narrowed as it is passed on.

Macaroons~\citep{macaroons2014} are the reference construction for attenuable bearer
credentials. A macaroon is a token whose holder can append \emph{caveats} that restrict
its use, producing a strictly less powerful token without contacting the issuer; the chain
of caveats is bound by a keyed-HMAC construction so that no caveat can be dropped. The
essential property for delegation is monotone attenuation: a derived credential can only
be a subset of the one it derives from.

agent-manifest~\citep{agentmanifest2026} implements this discipline as a signed A2A
delegation chain. Each hop is a \texttt{DelegationCredential} whose signed body names the
issuer, the subject, the granted scope, a depth, and the parent hop's identifier. The body
is signed with Ed25519~\citep{eddsa2017} over a canonical JSON encoding aligned with RFC
8785~\citep{jcs2020}, so the exact bytes that were signed are reconstructable by any
verifier. Verification enforces signature validity at every hop, continuity of the
issuer-to-subject and parent-to-child links, a bounded depth, non-repetition of credential
identifiers to defeat replay, and the subset check on scope that enforces attenuation. cA2A
implements the same semantics in its own verifier, which the runtime calls on every inbound
peer request, and signs over RFC 8785 canonical bytes so that credentials are
cross-verifiable with agent-manifest. The attenuation subset check is what forecloses the
confused-deputy escalation.

\subsection{Trusted Execution Environments and Remote Attestation}

A Trusted Execution Environment (TEE) is a CPU-enforced isolation boundary that provides
two properties cA2A depends on: \emph{confidentiality}, so that memory inside the TEE is
inaccessible to software outside it including the hypervisor and host OS, and
\emph{remote attestation}, so that the TEE can produce a hardware-signed measurement report
proving to a remote verifier which code is running inside. AMD SEV-SNP~\citep{amdsevsnp2023}
and Intel TDX~\citep{inteltdx2023} are the two mainstream VM-level implementations; a survey
of confidential machine learning provides broader context~\citep{confai2024}.

The measurement report is the primitive of interest. Both platforms let the enclave place
application-chosen bytes into a report field bound to the hardware measurement, which is how
a public key can be tied to a specific attested measurement. The IETF RATS
architecture~\citep{rats2021} gives the vendor-neutral roles for this exchange, an attester
producing evidence and a verifier appraising it, and the Entity Attestation Token
(EAT)~\citep{eat2023} gives a portable token format for carrying attestation claims. cA2A
binds a peer's channel key to its attested measurement through this mechanism, so that a
caller can seal a payload to a key the appraised evidence vouches for. Whether that key is in
fact confined to the measured code depends on the platform and on what the measurement
covers; Sections~\ref{sec:claims} and~\ref{sec:after} state what has and has not been shown.

\subsection{Cedar Policy Language}

Cedar~\citep{cedar2023} is a policy language and evaluation engine that expresses access
control as \texttt{permit} and \texttt{forbid} rules over a principal, action, and
resource, with optional conditions over a context object. Cedar's structural typing keeps
evaluation decidable in bounded time, which matters when a policy check sits on the
critical path of accepting a delegated task. cA2A uses Cedar to model the intersection of a
forwarded delegation scope with a peer's local policy, so that a peer never acts beyond
either the authority delegated to it or its own operator's rules. Property~\ref{prop:c3}
states this intersection and Section~\ref{sec:eval-c3} evaluates it.

\subsection{Provenance and Transparency}

An offline-verifiable record of who delegated what to whom is a provenance problem. The
W3C PROV data model~\citep{prov2013} gives the vocabulary for describing entities,
activities, and the agents responsible for them, and the relationships that link a
derived artifact to what it was derived from. cA2A's per-hop records form a delegation DAG
in exactly this shape: each record names its parent, so the set of records across a
multi-hop workflow reconstructs the delegation tree.

A transparency log makes a registered record non-repudiable and lets anyone check that it was
included; it cannot reveal a record that was never registered. Certificate
Transparency~\citep{ct2013} established the pattern of an append-only, cryptographically
verifiable log for security-critical statements, and the IETF SCITT
architecture~\citep{scitt2024} generalizes it to supply-chain claims, giving neutral roles
for registering and auditing signed statements. cA2A follows the same profile-over-standard
posture toward SCITT that it takes toward A2A: it binds to the transparency architecture
rather than reinventing it.

\subsection{TRACE as the Record Format}

TRACE (Trust, Runtime Attestation, and Compliance Evidence)~\citep{trace2025} is the open
record format that cA2A emits per delegation hop. A TRACE record is a signed JSON document
that carries a subject, a runtime measurement, and profile-specific context. cA2A defines
an A2A profile of TRACE that adds a delegation-link block to each record: the hash of the
parent hop's record and the identifier of the delegation credential the hop acted under. A
root hop omits the block; every subsequent hop references its parent. Because each record
is signed and references its parent by hash, a verifier can rebuild the parent links,
confirm each hop acted under a credential whose chain verifies, and confirm that each record
was signed by the key the chain names as that hop's delegate. These checks run offline, given
the recipient's trusted keys and the delegation chain. They establish what the signers
declared, not that the declared work completed or that the submitted path is the complete
history, and the current verifiers accept one ordered root-to-leaf path rather than an
arbitrary branching graph. This is the record substrate for the provenance property in
Section~\ref{sec:claims}.
